\phantomsection
\color{unidarkgreen}\section*{\centering{\large{ПЕРЕЧЕНЬ СОКРАЩЕНИЙ}}}
\addcontentsline{toc}{section}{Перечень сокращений}

% ===== НАСТРОЙКИ =====
\newlength{\abbrgap}
\setlength{\abbrgap}{1.5em}   % <-- расстояние между "--" и расшифровкой (меняйте здесь)

% Автоматическая ширина метки: самое длинное сокращение + тире + пробел + запас
\newlength{\abbrwidth}
\settowidth{\abbrwidth}{УПНКП}   % самое длинное сокращение
\addtolength{\abbrwidth}{3.5em}% запас под "-- " и отступ
% ======================

{\color{white}\fontsize{0.1pt}{0.1pt}\selectfont<begABBRS>}
\color{black}

\begin{list}{}%
{%
  \setlength{\labelwidth}{\abbrwidth}%
  \setlength{\labelsep}{0pt}%      <-- теперь весь отступ задаётся через \abbrgap в метке
  \setlength{\leftmargin}{\dimexpr\labelwidth+\labelsep\relax}%
  \setlength{\itemindent}{0pt}%
  \setlength{\parsep}{0pt}%
  \setlength{\itemsep}{0pt}
  \renewcommand{\makelabel}[1]{#1 --\hspace{\abbrgap}}%   <-- метка = сокращение + "--" + пробел
}
\item[IEC]  International Electrotechnical Commission;
\item[PPS] Pulse Per Second;
\item[АПВ] автоматическое повторное включение;
\item[АУВ] автоматика управления выключателем;
\item[БИ] блок испытательный;
\item[БСК] батарея статических конденсаторов;
\item[ДВ] дискретный вход;
\item[ЕЭС] единая энергетическая система;
\item[ЗНФ] защита от непереключения фаз;
\item[ИО] измерительный орган / избирательный орган;
\item[ИЧМ] интерфейс человек-машина;
\item[ИЭУ] интеллектуальное электронное устройство;
\item[КО] компенсационная обмотка;
\item[КОР] компенсационная обмотка реактора;
\item[ЛЭП] линия электропередачи;
\item[НКУ] низковольтное комплектное устройство;
\item[НПП] научно-производственное предприятие;
\item[НТД] нормативно-техническая документация;
\item[ОАПВ] однофазное автоматическое повторное включение;
\item[ОПУ] общеподстанционный пункт управления;
\item[ОТ] оперативный ток;
\item[ПА] противоаварийная автоматика;
\item[ПК] предохранительный клапан;
\item[ПО] пусковой орган / программное обеспечение;
\item[ПС] предупредительная сигнализация;
\item[ПУЭ] правила устройства электроустановок;
\item[РАС] регистратор аварийных событий;
\item[РЗА] релейная защита и автоматика;
\item[РПВ] реле положения «включено»;
\item[РПО] реле положения отключено;
\item[РФ] реле фиксации;
\item[РЭ] руководство по эксплуатации;
\item[СИ] средства измерений;
\item[СКРМ] средства компенсации реактивной мощности;
\item[ТК] токовый контроль / текущий контроль;
\item[ТН] трансформатор напряжения;
\item[ТО] токовая отсечка;
\item[ТТ] трансформатор тока;
\item[ТУ] телеуправление / телеускорение / технические условия;
\item[УПНКП] устройство принудительной несинхронной коммутации полюсов;
\item[УРОВ] устройство резервирования при отказе выключателя;
\item[УХЛ] умеренный и холодный климат;
\item[УШР] управляемый шунтирующий реактор;
\item[ШР] шунтирующий реактор;
\item[ШЭТ] шкаф электротехнический типовой;
\item[ЭМВ] электромагнит включения;
\item[ЭМО] электромагнит отключения.
{\color{white}\fontsize{0.1pt}{0.1pt}\selectfont<endABBRS>}
\end{list}
